\documentclass[11pt]{article}
\usepackage[margin=1in]{geometry}
\usepackage{amsmath,amssymb,amsthm,mathtools,hyperref,array}

\newtheorem{assumption}{Assumption}
\newtheorem{lemma}{Lemma}
\newtheorem{proposition}{Proposition}
\newtheorem{theorem}{Theorem}
\newtheorem{remark}{Remark}

\newcommand{\R}{\mathbb{R}}
\newcommand{\E}{\mathbb{E}}
\newcommand{\eps}{\varepsilon}

\title{Scalar And Finite-Dimensional Perron Theory Appendix For MESA}
\author{Working proof supplement}
\date{May 16, 2026}

\begin{document}
\maketitle

\section{Purpose}

This appendix gives a submission-style proof of the scalar amplification result
used by the MESA manuscript and a finite-dimensional Perron-visible
reduced-form bridge for multitype Hawkes order flow. It also verifies the
finite-state robust Bellman layer used by the numerical DP and a compact
continuous-intensity HJBI verification theorem for a truncated multitype
Hawkes control game. It also gives an untruncated Lyapunov bridge and weighted
comparison theorem under explicit robust subcriticality and weighted boundary
assumptions, plus a local differentiability theorem for interior robust quote
optimizers. It intentionally proves the narrow result that the current paper
can defend.

\section{Notation Checklist}

\begin{center}
\small
\begin{tabular}{>{\raggedright\arraybackslash}p{0.18\linewidth}>{\raggedright\arraybackslash}p{0.28\linewidth}>{\raggedright\arraybackslash}p{0.43\linewidth}}
\hline
Symbol & Domain & Meaning \\
\hline
$\rho$ & $[0,1)$ & Scalar branching ratio or Perron root
$\rho(\Gamma)$, depending on context. \\
$\Gamma$ & $\R_+^{d\times d}$ & Multitype Hawkes branching matrix. \\
$\mu,\lambda$ & $\R_+^d$ & Baseline intensity and current intensity vector. \\
$r,\ell$ & $\R_+^d$ & Right and left Perron vectors, normalized by
$\ell^\top r=1$. \\
$E,E_{\rm rel},E_{\rm abs}$ & Matrix perturbations & Generic, relative-slack,
and absolute branching-matrix perturbations. \\
$\eps,\eps_{\rm rel}$ & positive scalars & Absolute and relative ambiguity
radii. \\
$V_{\rm risk},V_d,V_{\rm bi}$ & real-valued functions & Reduced risk proxies
used to isolate the resolvent singularity. \\
$q$ & $\{-Q,\ldots,Q\}$ & Inventory state in the Bellman and HJBI games. \\
$K$ & compact subset of $\R_+^d$ & Truncated continuous-intensity domain. \\
$U,\mathcal A$ & compact or finite sets & Quote/no-quote controls and
rectangular ambiguity set. \\
$L^{u,\Gamma},\mathcal L^{u,\Gamma}$ & generators & Finite-state and
continuous-intensity controlled Hawkes generators. \\
$W_v,\kappa$ & Lyapunov function and drift margin & Weighted intensity
Lyapunov function and common subcriticality margin. \\
$J_i,K_R$ & jump map and compact set & Untruncated Hawkes mark-$i$ jump map
and Lyapunov sublevel truncation set $\{W_v\le R\}$. \\
$U,V,F$ & weighted viscosity functions & Generic subsolution, supersolution,
and solution candidate in the comparison theorem. \\
$\delta,\theta$ & quote vector and structural parameter & Local quote action
and finite-dimensional branching/ambiguity parameter in the optimizer
sensitivity theorem. \\
$\mathfrak H,\delta^*$ & reduced Hamiltonian and optimizer & Robust local
Hamiltonian after Nature's minimization and its interior quote maximizer. \\
$T,g,r$ & finite horizon, payoff, reward & Terminal time, terminal payoff, and
running reward. \\
\hline
\end{tabular}
\end{center}

All asymptotic constants below are uniform on the compact regularity classes
specified in the relevant theorem statements. The same letters may denote a
scalar object in the scalar section and its multitype analogue later; each
section restates the needed domain.

\section{Scalar Hawkes Setup}

Let $N_t$ be a stationary linear Hawkes process with exponential kernel
\[
\lambda_t = \mu + \int_0^{t^-}\beta \rho e^{-\beta(t-s)}\,dN_s,
\qquad
0\le \rho < 1,\quad \mu>0,\quad \beta>0.
\]
Here $\rho$ is the branching ratio. The stationary mean intensity is
\[
\bar\lambda(\rho)=\frac{\mu}{1-\rho}.
\]

We use a reduced market-making risk proxy
\[
V_{\mathrm{risk}}(\rho)=C(1-\rho)^{-a},
\qquad C>0,
\]
where $a=2$ is the normalized-throughput variance convention used for the
headline square-law result. If stationary throughput is allowed to enter the
risk term separately, then $a$ may be larger; the proofs below make the exponent
explicit.

\begin{assumption}[Local quote sensitivity]
For inventory $q$ in a compact grid and quote controls in a compact interval,
the robust half-spread correction is locally linear in the worst-case increase
of $V_{\mathrm{risk}}$:
\[
\psi_\eps(\rho)
=K\left[\sup_{\tilde\rho\in\mathcal U_\eps(\rho)}
V_{\mathrm{risk}}(\tilde\rho)-V_{\mathrm{risk}}(\rho)\right]
+o\left(\sup_{\tilde\rho\in\mathcal U_\eps(\rho)}
|V_{\mathrm{risk}}(\tilde\rho)-V_{\mathrm{risk}}(\rho)|\right),
\]
with $K>0$ independent of $\rho$ near one.
\end{assumption}

This assumption is the scalar version of differentiating the first-order
condition for an interior quote optimizer. The optimizer-sensitivity theorem
below gives sufficient conditions under which it follows from a uniformly
negative second derivative in quote space and a smooth active worst-case
selector for the robust Hamiltonian.

\section{Resolvent And Variance Scaling}

\begin{lemma}[Scalar resolvent]
For scalar Hawkes branching ratio $\rho<1$,
\[
(1-\rho)^{-1}
\]
is the scalar Hawkes resolvent. Its square scales as $(1-\rho)^{-2}$.
\end{lemma}

\begin{proof}
The stationary mean equation is
$\bar\lambda=\mu+\rho\bar\lambda$. Therefore
$\bar\lambda=(1-\rho)^{-1}\mu$. The linear response of total descendants of one
immigrant is the geometric series $\sum_{n\ge0}\rho^n=(1-\rho)^{-1}$, proving
the resolvent identity. Squaring the leading response gives
$(1-\rho)^{-2}$.
\end{proof}

\begin{lemma}[Derivative of the risk proxy]
Let $V_{\mathrm{risk}}(\rho)=C(1-\rho)^{-a}$ for $a>0$. Then
\[
\frac{d}{d\rho}V_{\mathrm{risk}}(\rho)
=aC(1-\rho)^{-(a+1)}.
\]
\end{lemma}

\begin{proof}
Direct differentiation gives
$d(1-\rho)^{-a}/d\rho=a(1-\rho)^{-(a+1)}$.
\end{proof}

\section{Ambiguity Normalizations}

We distinguish two ambiguity sets.

\paragraph{Relative critical-slack ambiguity.}
\[
\mathcal U^{\mathrm{rel}}_\eps(\rho)
=\{\tilde\rho:\ 0\le \tilde\rho\le \rho+\eps(1-\rho),\ \tilde\rho<1\}.
\]

\paragraph{Absolute branching-ratio ambiguity.}
\[
\mathcal U^{\mathrm{abs}}_\eps(\rho)
=\{\tilde\rho:\ 0\le \tilde\rho\le \rho+\eps,\ \tilde\rho<1\}.
\]

\begin{theorem}[Relative-slack amplification]
Fix $a>0$ and $0<\eps<1$. Under relative critical-slack ambiguity,
\[
\sup_{\tilde\rho\in\mathcal U^{\mathrm{rel}}_\eps(\rho)}
V_{\mathrm{risk}}(\tilde\rho)-V_{\mathrm{risk}}(\rho)
=C\left((1-\eps)^{-a}-1\right)(1-\rho)^{-a}.
\]
Consequently, for small $\eps$,
\[
\psi_\eps(\rho)
=KaC\,\eps(1-\rho)^{-a}+O(\eps^2(1-\rho)^{-a})+o(\eps(1-\rho)^{-a}).
\]
In the normalized-throughput convention $a=2$, the premium has the square law
$\psi_\eps(\rho)=\Theta(\eps(1-\rho)^{-2})$.
\end{theorem}

\begin{proof}
The risk proxy is increasing in $\rho$, so Nature chooses
$\tilde\rho^\star=\rho+\eps(1-\rho)$. Then
\[
1-\tilde\rho^\star = 1-\rho-\eps(1-\rho)=(1-\eps)(1-\rho).
\]
Substitution gives
\[
V_{\mathrm{risk}}(\tilde\rho^\star)-V_{\mathrm{risk}}(\rho)
=C(1-\rho)^{-a}\left((1-\eps)^{-a}-1\right).
\]
The Taylor expansion
$(1-\eps)^{-a}=1+a\eps+O(\eps^2)$ proves the small-$\eps$ expression. The
quote-premium statement follows from local quote sensitivity. Since all
constants are positive and independent of $\rho$ near one, the result is
$\Theta(\eps(1-\rho)^{-a})$.
\end{proof}

\begin{theorem}[Absolute-ambiguity derivative law]
Fix $a>0$ and suppose $\eps=o(1-\rho)$. Under absolute ambiguity,
\[
\sup_{\tilde\rho\in\mathcal U^{\mathrm{abs}}_\eps(\rho)}
V_{\mathrm{risk}}(\tilde\rho)-V_{\mathrm{risk}}(\rho)
=aC\,\eps(1-\rho)^{-(a+1)}
+O(\eps^2(1-\rho)^{-(a+2)}).
\]
In the normalized-throughput convention $a=2$, the premium scales as
$\Theta(\eps(1-\rho)^{-3})$.
\end{theorem}

\begin{proof}
Again Nature chooses $\tilde\rho^\star=\rho+\eps$ when this remains below one.
Taylor expanding around $\rho$,
\[
V_{\mathrm{risk}}(\rho+\eps)-V_{\mathrm{risk}}(\rho)
=\eps V'_{\mathrm{risk}}(\rho)
+\frac{\eps^2}{2}V''_{\mathrm{risk}}(\xi)
\]
for some $\xi\in(\rho,\rho+\eps)$. By the derivative lemma,
$V'_{\mathrm{risk}}(\rho)=aC(1-\rho)^{-(a+1)}$. Also
$V''_{\mathrm{risk}}(\xi)=a(a+1)C(1-\xi)^{-(a+2)}$. The condition
$\eps=o(1-\rho)$ implies $1-\xi$ is comparable to $1-\rho$, giving the stated
remainder. Local quote sensitivity transfers the scaling to $\psi_\eps$.
\end{proof}

\section{Multitype Perron Extension}

\begin{proposition}[First-order Perron perturbation]
Let $\Gamma$ be nonnegative and irreducible with simple Perron root
$\rho<1$, normalized left/right Perron vectors $\ell^\top r=1$, and spectral
gap bounded below. Then for a perturbation $E$,
\[
\rho(\Gamma+E)-\rho(\Gamma)=\ell^\top E r+O(\|E\|_F^2).
\]
The Frobenius-unit perturbation maximizing the first-order increase is
$E^\star = \ell r^\top/\|\ell r^\top\|_F$.
\end{proposition}

\begin{proof}
The first statement is the standard analytic perturbation expansion for a
simple eigenvalue. The derivative of the Perron root in direction $E$ is
$\ell^\top E r=\langle E,\ell r^\top\rangle_F$. Under a Frobenius norm
constraint, Cauchy--Schwarz gives
\[
\langle E,\ell r^\top\rangle_F
\le \|E\|_F\|\ell r^\top\|_F,
\]
with equality when $E$ is aligned with $\ell r^\top$.
\end{proof}

\begin{remark}
This proposition justifies the Perron-aligned adversary used in the experiments
when the leading eigenvalue is simple and well separated. It also explains why
near-degenerate spectra are an ablation, not a nuisance parameter: the lower
bound can weaken when Perron regularity fails.
\end{remark}

\begin{theorem}[Finite-dimensional Perron-visible ambiguity premium]
Let $\Gamma\in\R_+^{d\times d}$ be irreducible, primitive, and subcritical
with simple Perron root $\rho<1$. Assume a uniform spectral gap
\[
\rho-\max_{j>1}|\lambda_j(\Gamma)|\ge g_0>0
\]
and uniformly bounded Perron eigenprojection conditioning. Let $r,\ell>0$ be
normalized by $\ell^\top r=1$. For fixed $\mu,w\in\R_+^d$, define
\[
V_d(\Gamma)=C\left(w^\top(I-\Gamma)^{-1}\mu\right)^2,
\qquad C>0,
\]
and assume Perron visibility
\[
A(\Gamma):=(w^\top r)(\ell^\top\mu)\ge a_0>0.
\]
Then
\[
V_d(\Gamma)=C A(\Gamma)^2(1-\rho)^{-2}+O((1-\rho)^{-1}).
\]

If a feasible Perron-aligned relative-slack perturbation satisfies
\[
\rho(\Gamma+E_{\rm rel})
=\rho+\eps_{\rm rel}(1-\rho)+O(\eps_{\rm rel}^2(1-\rho)),
\]
then
\[
V_d(\Gamma+E_{\rm rel})-V_d(\Gamma)
=\Theta\!\left(\eps_{\rm rel}(1-\rho)^{-2}\right).
\]

If a feasible Perron-aligned absolute perturbation satisfies
\[
\rho(\Gamma+E_{\rm abs})=\rho+\eps+O(\eps^2),
\qquad \eps=o(1-\rho),
\]
then
\[
V_d(\Gamma+E_{\rm abs})-V_d(\Gamma)
=\Theta\!\left(\eps(1-\rho)^{-3}\right).
\]
Under the local quote-sensitivity assumption, the robust half-spread premium
has the same respective orders.
\end{theorem}

\begin{proof}
By the Riesz projection formula for the isolated Perron eigenvalue,
\[
(I-\Gamma)^{-1}=\frac{r\ell^\top}{1-\rho}+R,
\qquad \|R\|\le C(g_0),
\]
where the constant is uniform on the Perron-regular class. Therefore
\[
w^\top(I-\Gamma)^{-1}\mu
=\frac{(w^\top r)(\ell^\top\mu)}{1-\rho}+O(1).
\]
Squaring gives the displayed expansion for $V_d$ because Perron visibility
keeps the leading coefficient bounded away from zero.

For perturbations, the first-order Perron expansion gives
\[
\rho(\Gamma+E)-\rho(\Gamma)
=\ell^\top E r+O(\|E\|_F^2/g_0).
\]
The Perron eigenvectors and the visibility coefficient vary by lower-order
terms under the stated local perturbations. Expanding
$(1-\rho-\Delta\rho)^{-2}-(1-\rho)^{-2}$ gives the relative-slack order when
$\Delta\rho=\eps_{\rm rel}(1-\rho)+O(\eps_{\rm rel}^2(1-\rho))$, and the
absolute-ambiguity derivative law when $\Delta\rho=\eps+O(\eps^2)$ with
$\eps=o(1-\rho)$. The quote-premium statement follows from local quote
sensitivity.
\end{proof}

\section{Finite-State Robust Bellman Bridge}

This section closes one control-theoretic gap for the implemented experiment
package. It proves the robust dynamic-programming equation for the finite
inventory, finite-regime approximation used numerically. By itself, this
finite-state theorem is not the continuous-intensity multitype Hawkes HJBI
bridge; that compact bridge is stated in the next section.

Let $\mathcal X$ be a finite state space whose elements encode inventory,
cash-normalized wealth state, and a finite Hawkes-regime or intensity bin. Let
$\mathcal U(x)$ be a nonempty finite set of admissible quote/no-quote actions
at state $x$, and let $\mathcal A=\{\Gamma_1,\ldots,\Gamma_M\}$ be a finite
rectangular ambiguity set of stable branching matrices. For each
$(u,\Gamma)\in\mathcal U(x)\times\mathcal A$, let
$L^{u,\Gamma}(x,y)$ be a conservative generator on $\mathcal X$:
\[
L^{u,\Gamma}(x,y)\ge0\quad (y\ne x),
\qquad
L^{u,\Gamma}(x,x)=-\sum_{y\ne x}L^{u,\Gamma}(x,y).
\]
Assume bounded running rewards $r(x,u,\Gamma)$ and bounded terminal payoff
$g(x)$.

\begin{theorem}[Finite-state robust Hawkes-regime Bellman verification]
For horizon $T<\infty$, the rectangular robust market-making game
\[
\sup_{u}\inf_{\Gamma}\E\left[
\int_t^T r(X_s,u_s,\Gamma_s)\,ds+g(X_T)
\,\middle|\,X_t=x\right],
\]
where $u_s$ and $\Gamma_s$ are predictable Markov selectors from
$\mathcal U(X_s)$ and $\mathcal A$, has a value $V(t,x)$. The value is the
unique bounded solution of the finite HJBI system
\[
-\partial_t V(t,x)
=\max_{u\in\mathcal U(x)}\min_{\Gamma\in\mathcal A}
\left[
r(x,u,\Gamma)+
\sum_{y\in\mathcal X}L^{u,\Gamma}(x,y)V(t,y)
\right],
\qquad
V(T,x)=g(x).
\]
Moreover, deterministic Markov selectors attaining the displayed finite
maximum and minimum are optimal. The same statement holds for an explicit
Euler time grid, where the ODE is replaced by the backward Bellman recursion.
\end{theorem}

\begin{proof}
Because $\mathcal X$, $\mathcal U(x)$, and $\mathcal A$ are finite and rewards
are bounded, the Hamiltonian
\[
H_x(v)=\max_{u\in\mathcal U(x)}\min_{\Gamma\in\mathcal A}
\left[r(x,u,\Gamma)+\sum_y L^{u,\Gamma}(x,y)v(y)\right]
\]
is globally Lipschitz in $v\in\R^{|\mathcal X|}$. Indeed, for any $v,w$,
the difference between two generator pairings is bounded by
$2\Lambda\|v-w\|_\infty$, where
$\Lambda=\max_{x,u,\Gamma}\sum_{y\ne x}L^{u,\Gamma}(x,y)<\infty$; taking
finite maxima and minima preserves this Lipschitz constant. Hence the terminal
value problem $-\dot V=H(V)$, $V(T)=g$, has a unique bounded classical solution
by the Picard--Lindelof theorem for finite-dimensional ODEs.

For a small step $\Delta t$, conditioning on the first jump of the controlled
Markov chain gives
\[
V(t,x)=
\max_u\min_\Gamma
\left[
r(x,u,\Gamma)\Delta t+
\sum_y P_{\Delta t}^{u,\Gamma}(x,y)V(t+\Delta t,y)
\right]+o(\Delta t),
\]
where $P_{\Delta t}^{u,\Gamma}=I+\Delta t L^{u,\Gamma}+o(\Delta t)$. Subtract
$V(t+\Delta t,x)$, divide by $\Delta t$, and let $\Delta t\downarrow0$ to
recover the HJBI system. Conversely, applying Dynkin's formula to the unique
solution under arbitrary Markov selectors and using the pointwise
max-min inequality shows that the solution dominates every admissible
controller against worst-case Nature and is achieved by any selectors attaining
the finite maximum and minimum. Such selectors exist because the action and
ambiguity sets are finite. The discrete-time statement follows by the same
one-step conditioning argument without taking the limit.
\end{proof}

\begin{remark}
This theorem justifies the finite-scenario Bellman layer used in the code. If
the scenario set contains Perron-aligned relative-slack perturbations, the
finite game can express the robust quote/no-quote response to those regimes.
The theorem does not prove the continuum multitype Hawkes HJBI or
differentiability of the true optimizer with respect to $\Gamma$. The next
section gives the compact continuous-intensity bridge.
\end{remark}

\section{Compact Continuous-Intensity HJBI Bridge}

This section is the compact-domain layer of the multitype control proof. It
states the compact piecewise-deterministic Markov-process HJBI theorem under
the usual dynamic-programming and comparison hypotheses for the admissible
control class. The next two sections remove the projection through a
Lyapunov/tightness argument and a weighted comparison theorem, and the
consolidated theorem below records the resulting untruncated statement.

Let the state be $X_t=(q_t,\lambda_t)$, where
$q_t\in\{-Q,\ldots,Q\}$ is inventory and
$\lambda_t\in K\subset\R_+^d$ is the vector of Hawkes intensities. The compact
set $K$ is interpreted as the numerical truncation domain. Between jumps,
\[
d\lambda_t=\beta(\mu-\lambda_t)\,dt.
\]
At a mark-$i$ jump, inventory changes by a bounded increment $\zeta_i(q,u)$
and intensity is updated by a truncated Hawkes jump map
\[
\lambda_{t}=\Pi_K\left(\lambda_{t^-}+\beta\Gamma e_i\right),
\]
where $\Pi_K$ is a fixed Lipschitz projection into $K$. The market maker chooses
$u\in U(q,\lambda)$ and Nature chooses
$\Gamma\in\mathcal A$, where $U(q,\lambda)$ and $\mathcal A$ are compact. Let
$\Lambda_i(q,\lambda,u,\Gamma)$ be the controlled mark intensities after quote
placement and censoring by inventory limits. Define the generator on smooth
test functions $\varphi$ by
\[
\begin{aligned}
\mathcal L^{u,\Gamma}\varphi(q,\lambda)
={}&\beta(\mu-\lambda)\cdot\nabla_\lambda\varphi(q,\lambda)\\
&+\sum_{i=1}^d \Lambda_i(q,\lambda,u,\Gamma)
\left[
\varphi(q+\zeta_i(q,u),\Pi_K(\lambda+\beta\Gamma e_i))
-\varphi(q,\lambda)
\right].
\end{aligned}
\]

\begin{assumption}[Compact Hawkes control regularity]
The sets $K$, $\mathcal A$, and $U(q,\lambda)$ are compact, with
$U(q,\lambda)$ having a closed graph. The ambiguity set satisfies
$\sup_{\Gamma\in\mathcal A}\rho(\Gamma)\le\rho_{\max}<1$. The functions
$r(q,\lambda,u,\Gamma)$, $g(q,\lambda)$, and
$\Lambda_i(q,\lambda,u,\Gamma)$ are bounded and continuous, and are uniformly
Lipschitz in $(q,\lambda)$ on each inventory slice. The jump map above remains
inside $K$, and the total jump rate is uniformly bounded by $\bar\Lambda$.
Admissible controllers are nonanticipative predictable controls; Nature's
choice of $\Gamma$ is rectangular and predictable over time. The compact PDMP
game satisfies the dynamic-programming principle for this admissible class, and
bounded viscosity sub/supersolutions of the compact HJBI satisfy the standard
comparison principle.
\end{assumption}

\begin{theorem}[Compact continuous-intensity Hawkes HJBI verification]
Under compact Hawkes control regularity, the robust finite-horizon rectangular
game
\[
V(t,q,\lambda)=
\sup_u\inf_\Gamma
\E\left[
\int_t^T r(q_s,\lambda_s,u_s,\Gamma_s)\,ds+
g(q_T,\lambda_T)
\,\middle|\,q_t=q,\lambda_t=\lambda
\right]
\]
has a bounded uniformly continuous value function. It is the unique bounded
viscosity solution, for the compact admissible class specified above, on
$[0,T]\times\{-Q,\ldots,Q\}\times K$ of
\[
-\partial_t V=
\sup_{u\in U(q,\lambda)}\inf_{\Gamma\in\mathcal A}
\left\{
r(q,\lambda,u,\Gamma)+
\mathcal L^{u,\Gamma}V(q,\lambda)
\right\},
\qquad
V(T,q,\lambda)=g(q,\lambda).
\]
If, in addition, the value is $C^1$ in $(t,\lambda)$ on each inventory slice
and the displayed supremum and infimum admit measurable selectors
$u^\star(t,q,\lambda)$ and $\Gamma^\star(t,q,\lambda)$ satisfying Isaacs'
pointwise saddle condition, then these selectors verify the game value. In the
absence of exact selectors, measurable $\eta$-optimizers are $\eta$-optimal.
\end{theorem}

\begin{proof}
The assumptions define a controlled piecewise-deterministic Markov process on
the compact state space
$\{-Q,\ldots,Q\}\times K$. Bounded rates rule out explosion, and compactness
with continuity gives measurable selectors for the relaxed one-step
Hamiltonians. The assumed compact PDMP dynamic-programming principle gives
\[
V(t,x)=
\sup_u\inf_\Gamma
\E\left[
\int_t^{t+h} r(X_s,u_s,\Gamma_s)\,ds+
V(t+h,X_{t+h})\middle|X_t=x
\right]
\]
for $x=(q,\lambda)$.

Let $\phi$ be a smooth test function touching $V$ from above at
$(t,x)$. Apply the dynamic programming principle with nearly optimal controls
over $[t,t+h]$, use the first-jump expansion for the PDMP, subtract
$\phi(t,x)$, divide by $h$, and let $h\downarrow0$. This gives the viscosity
subsolution inequality. The same argument with a test function touching from
below, freezing any admissible control and allowing Nature to use a nearly
minimizing response, gives the supersolution inequality.

For comparison, the Hamiltonian is proper and uniformly Lipschitz in the value
argument because total jump rates are bounded by $\bar\Lambda$, rewards are
bounded, and coefficients are uniformly Lipschitz on the compact domain. The
assumed compact comparison principle therefore implies that any bounded
upper-semicontinuous subsolution is dominated by any bounded
lower-semicontinuous supersolution with ordered terminal data. Perron's method
and the dynamic programming principle give existence, hence uniqueness of the
bounded viscosity solution.

If $V$ is classically differentiable in $(t,\lambda)$, applying Dynkin's
formula to $V(s,X_s)$ under arbitrary selectors and using the pointwise
max-min inequalities proves that $V$ bounds the payoff from above and below.
Selectors attaining the saddle Hamiltonian make those inequalities equalities,
so they verify the value. Compactness gives measurable $\eta$-optimizers when
exact selectors are unavailable, yielding $\eta$-optimal controls.
\end{proof}

\begin{remark}
This theorem closes the compact truncated bridge conditional on the standard
compact-PDMP dynamic-programming and comparison hypotheses stated above. The
next theorem removes the compact intensity projection for stability and
truncation tails. The weighted comparison section then gives the untruncated
viscosity identification for pairs with a fixed weighted boundary condition at
infinity, and the optimizer-sensitivity section isolates no-quote/boundary
active-set differentiability as a separate nonsmooth control issue.
\end{remark}

\section{Untruncated Lyapunov Bridge}

This section removes the artificial compact intensity projection at the level
of nonexplosion, moment control, and truncation error. It is not a universal
comparison theorem for every unbounded HJBI; it supplies the missing
Lyapunov/tightness component under an explicit common subcriticality condition.

Let $\lambda_t\in\R_+^d$ evolve without projection:
\[
d\lambda_t=\beta(\mu-\lambda_t)\,dt,
\qquad
\lambda_t\mapsto \lambda_t+\beta\Gamma e_i
\quad\hbox{at a mark-}i\hbox{ jump.}
\]
Controls may thin or censor mark intensities, but not amplify them beyond the
Hawkes state: assume
\[
0\le \Lambda_i(q,\lambda,u,\Gamma)\le \lambda_i
\]
for all admissible $(q,\lambda,u,\Gamma)$. This covers passive quote refusal
and inventory censoring, which can reduce fill intensities but should not
create order-flow arrivals ex nihilo.

\begin{assumption}[Common Lyapunov subcriticality and tail integrability]
There exist $v\in\R_+^d$ with strictly positive entries and $\kappa\in(0,1]$
such that
\[
\Gamma^\top v\le (1-\kappa)v
\]
componentwise for every $\Gamma\in\mathcal A$. Rewards and terminal payoffs
have at most linear Lyapunov growth:
\[
|r(q,\lambda,u,\Gamma)|+|g(q,\lambda)|\le C(1+v^\top\lambda)
\]
uniformly over the finite inventory grid and admissible controls. Fix some
$p>1$. Since $\mathcal A$ is compact, the Hawkes jumps have uniformly bounded
$W_v$-increments, so the $p$th Lyapunov moment calculation below is finite
uniformly over $\mathcal A$.
\end{assumption}

\begin{theorem}[Untruncated Hawkes PDMP Lyapunov stability]
Under common Lyapunov subcriticality, the unprojected controlled Hawkes PDMP is
nonexplosive for every predictable admissible controller and Nature selector.
For
\[
W_v(\lambda)=1+v^\top\lambda,
\]
its generator satisfies the uniform drift bound
\[
\mathcal L^{u,\Gamma}W_v(q,\lambda)
\le
\beta v^\top\mu-\beta\kappa\, v^\top\lambda
\le C_0-C_1 W_v(\lambda)
\]
for constants $C_0,C_1>0$ depending only on
$(\beta,\mu,v,\kappa)$. Consequently, for every finite horizon $T$ and the
$p>1$ in the assumption,
\[
\sup_{u,\Gamma}\E_{q,\lambda}\left[\sup_{0\le s\le T} W_v(\lambda_s)^p\right]
\le C_{T,p} W_v(\lambda)^p,
\]
and the finite-horizon value is finite in the weighted class
$
|V(t,q,\lambda)|\le C_T W_v(\lambda).
$
Moreover, if $\tau_R=\inf\{t:W_v(\lambda_t)\ge R\}$, then
\[
\sup_{u,\Gamma}\mathbb P_{q,\lambda}(\tau_R\le T)
\le \frac{C_{T,p} W_v(\lambda)^p}{R^p}.
\]
Thus compact truncations on $\{W_v\le R\}$ approximate the untruncated game in
probability. For linearly growing payoffs, the exit-tail contribution is
uniformly controlled by
\[
\sup_{u,\Gamma}
\E_{q,\lambda}\!\left[
\sup_{s\le T}W_v(\lambda_s)\mathbf 1_{\{\tau_R\le T\}}
\right]
\le
\frac{C_{T,p}W_v(\lambda)^p}{R^{p-1}},
\]
so the compact-truncation payoff error vanishes locally uniformly as
$R\uparrow\infty$.
\end{theorem}

\begin{proof}
For the affine test function $W_v$, the deterministic drift contributes
\[
\beta(\mu-\lambda)\cdot\nabla_\lambda W_v
=\beta v^\top\mu-\beta v^\top\lambda.
\]
At a mark-$i$ jump, $W_v$ increases by
$\beta v^\top\Gamma e_i$. Therefore
\[
\mathcal L^{u,\Gamma}W_v
=\beta v^\top\mu-\beta v^\top\lambda
+\sum_{i=1}^d \Lambda_i(q,\lambda,u,\Gamma)\,\beta v^\top\Gamma e_i.
\]
Using $\Lambda_i\le\lambda_i$ and the common Lyapunov inequality gives
\[
\sum_i \Lambda_i v^\top\Gamma e_i
\le
\sum_i \lambda_i(\Gamma^\top v)_i
\le
(1-\kappa)\sum_i \lambda_i v_i.
\]
Substitution yields
\[
\mathcal L^{u,\Gamma}W_v
\le
\beta v^\top\mu-\beta\kappa v^\top\lambda.
\]
Since $W_v=1+v^\top\lambda$, this is the stated Foster--Lyapunov drift after
absorbing constants.

Dynkin's formula applied to $W_v(\lambda_{t\wedge\tau_R})$ gives a uniform
first-moment bound up to $\tau_R$. Gronwall's inequality gives
$\E W_v(\lambda_{t\wedge\tau_R})\le C_TW_v(\lambda_0)$ with constants
independent of $R$, $u$, and $\Gamma$. Letting $R\uparrow\infty$ proves
nonexplosion and the first-moment estimate.

For the $p$th moment, apply the generator to $W_v^p$. Each mark-$i$ jump adds
the bounded amount $c_i(\Gamma)=\beta v^\top\Gamma e_i$ to $W_v$. Taylor's
formula gives, for $W_v\ge1$,
\[
(W_v+c_i)^p-W_v^p\le p c_i W_v^{p-1}+C_p W_v^{p-2}.
\]
The deterministic drift contributes
$pW_v^{p-1}\beta(v^\top\mu-v^\top\lambda)$, while
\[
\sum_i\lambda_i c_i(\Gamma)\le \beta(1-\kappa)v^\top\lambda.
\]
Thus, uniformly over controls and Nature,
\[
\mathcal L^{u,\Gamma}W_v^p\le C_p W_v^{p-1}-p\beta\kappa W_v^p
\le C'_p-c_p W_v^p
\]
after increasing constants. Dynkin's formula and Gronwall yield
$\E W_v(\lambda_{t\wedge\tau_R})^p\le C_{T,p}W_v(\lambda_0)^p$ uniformly in
$R$, $u$, and $\Gamma$.

For the maximal estimate, write the
stopped semimartingale decomposition
\[
W_v(\lambda_{t\wedge\tau_R})
=W_v(\lambda_0)+A_{t\wedge\tau_R}+M_{t\wedge\tau_R}.
\]
The drift variation satisfies
$|A_{t\wedge\tau_R}|\le C\int_0^{t\wedge\tau_R}W_v(\lambda_s)\,ds$. The
predictable quadratic variation of $M$ is also bounded by
$C\int_0^{t\wedge\tau_R}W_v(\lambda_s)\,ds$, because each mark-$i$ jump changes
$W_v$ by the fixed amount $\beta v^\top\Gamma e_i$ and the total marked jump
rate is at most $\sum_i\lambda_i$. Burkholder--Davis--Gundy with exponent $p$
and Gronwall then
give
\[
\E\sup_{s\le T}W_v(\lambda_{s\wedge\tau_R})^p\le C_{T,p}W_v(\lambda_0)^p,
\]
uniformly in $R$. Letting $R\uparrow\infty$ proves the displayed maximal
moment bound.
Markov's inequality then gives the exit bound for $\tau_R$.

The value-growth statement follows by integrating the assumed linear growth of
running rewards and terminal payoff against the moment bound. On
$\{W_v\le R\}$, the compact HJBI theorem applies to the truncated game. The
$p$th maximal moment and Holder's inequality give
\[
\E\left[
\sup_{s\le T}W_v(\lambda_s)\mathbf 1_{\{\tau_R\le T\}}
\right]
\le
\left(\E\sup_{s\le T}W_v(\lambda_s)^p\right)^{1/p}
\mathbb P(\tau_R\le T)^{1-1/p},
\]
which is bounded by $C_{T,p}W_v(\lambda_0)^p/R^{p-1}$. This controls the payoff
contribution from paths that leave the truncation set before $T$, proving
convergence in probability and tail-controlled approximation for linearly
growing payoffs.
\end{proof}

\begin{remark}
If $\mathcal A$ is finite and each matrix is strictly subcritical, the common
vector condition is easy to check by linear programming. For a general compact
ambiguity set, it is a substantive robust-stability assumption, stronger than
pointwise $\rho(\Gamma)<1$. This is the correct price of a uniform
untruncated theorem.
\end{remark}

\begin{remark}
Combining the compact HJBI theorem with the Lyapunov bridge gives an
untruncated value approximation theorem in the weighted class: solve on
$\{W_v\le R\}$, stop at the boundary, and let $R\uparrow\infty$. The next
section adds the corresponding viscosity comparison theorem in the
$W_v$-weighted class with an explicit boundary condition at infinity.
\end{remark}

\section{Weighted Comparison And Global Untruncated HJBI}

This section states the sufficient comparison condition that closes the
unbounded-domain HJBI bridge. It is deliberately explicit: without a weighted
boundary condition at infinity, unbounded intensity states can admit spurious
solutions.

\begin{assumption}[Weighted comparison regularity]
In addition to common Lyapunov subcriticality, assume:
\begin{enumerate}
\item The control graph $U(q,\lambda)$ and ambiguity set $\mathcal A$ are
compact, and measurable $\eta$-selectors exist for the local Hamiltonian.
\item The controlled rates $\Lambda_i(q,\lambda,u,\Gamma)$ are continuous,
bounded above by $\lambda_i$, and locally Lipschitz in $\lambda$, uniformly on
compact subsets and uniformly over $(u,\Gamma)$.
\item The jump maps
\[
J_i(q,\lambda,u,\Gamma)=
(q+\zeta_i(q,u),\,\lambda+\beta\Gamma e_i)
\]
are Lipschitz in $\lambda$ and satisfy
\[
W_v(J_i(q,\lambda,u,\Gamma))\le C_J W_v(\lambda).
\]
\item Rewards and terminal payoffs are continuous and have at most linear
$W_v$-growth.
\item Subsolutions $U$ and supersolutions $V$ are compared in the weighted
class
\[
|U(t,q,\lambda)|+|V(t,q,\lambda)|\le C W_v(\lambda)
\]
and satisfy the boundary condition at infinity
\[
\limsup_{R\to\infty}\;
\sup_{\substack{t\in[0,T],\,q\\ W_v(\lambda)\ge R}}
\frac{U(t,q,\lambda)-V(t,q,\lambda)}{W_v(\lambda)}
\le 0.
\]
\end{enumerate}
\end{assumption}

\begin{theorem}[Weighted comparison for the untruncated Hawkes HJBI]
Under weighted comparison regularity, let $U$ be an upper-semicontinuous
viscosity subsolution and $V$ a lower-semicontinuous viscosity supersolution
of the untruncated robust Hawkes HJBI
\[
-\partial_t F=
\sup_{u\in U(q,\lambda)}\inf_{\Gamma\in\mathcal A}
\left\{
r(q,\lambda,u,\Gamma)+\mathcal L^{u,\Gamma}F(q,\lambda)
\right\}
\]
on $[0,T)\times\{-Q,\ldots,Q\}\times\R_+^d$. If
$U(T,q,\lambda)\le V(T,q,\lambda)$ for all $(q,\lambda)$ and the weighted
boundary condition above holds, then
\[
U(t,q,\lambda)\le V(t,q,\lambda)
\]
for every $(t,q,\lambda)$. Consequently, the untruncated robust Hawkes control
problem has at most one viscosity solution in any subclass of the weighted
class with a fixed boundary condition at infinity; equivalently, comparison is
pairwise and requires the displayed weighted boundary ordering. The value
constructed by compact truncation is the unique solution in that subclass, and
the truncated solutions converge locally uniformly to it as $R\uparrow\infty$.
\end{theorem}

\begin{proof}
Fix $\delta>0$ and define the strict supersolution
\[
V^\delta(t,q,\lambda)
=V(t,q,\lambda)+\delta e^{a(T-t)}W_v(\lambda),
\]
with $a$ chosen large enough that the Lyapunov drift dominates the linear
growth constants in the Hamiltonian. The Foster--Lyapunov inequality gives
\[
-\partial_t\{\delta e^{a(T-t)}W_v\}
-\sup_u\inf_\Gamma \mathcal L^{u,\Gamma}
\{\delta e^{a(T-t)}W_v\}
\ge c\delta W_v(\lambda)-C\delta
\]
for some $c>0$ after increasing $a$. Thus $V^\delta$ is a strict
supersolution outside a fixed compact set, and the weighted boundary condition
implies that any positive maximum of $U-V^\delta$ occurs on
$[0,T]\times\{-Q,\ldots,Q\}\times K_R$ for some compact Lyapunov sublevel set
$K_R=\{W_v\le R\}$.

On this compact set, apply the standard doubling-of-variables argument to
\[
U(t,q,\lambda)-V^\delta(s,q,\lambda')
-\frac{|t-s|^2}{2\alpha}
-\frac{|\lambda-\lambda'|^2}{2\varepsilon}.
\]
The finite inventory coordinate is discrete, so at a positive maximum it can
be matched by taking the same $q$; otherwise the maximum over finitely many
inventory states is handled componentwise. The viscosity inequalities at the
maximizers, together with compactness of $U(q,\lambda)$ and $\mathcal A$, give
nearly maximizing/minimizing controls and ambiguity selectors. Continuity and
local Lipschitz regularity of rates, rewards, and jump maps control the drift
and jump terms. The nonlocal terms are handled by the viscosity nonlocal
doubling inequality:
\[
U(t,J_i(q,\lambda,u,\Gamma))-V^\delta(s,J_i(q,\lambda',u,\Gamma))
\]
is bounded by the same doubled maximum plus an
$O(|\lambda-\lambda'|^2/\varepsilon)+o(1)$ error, because each $J_i$ is
Lipschitz on $K_R$. Letting $\alpha,\varepsilon\downarrow0$ yields a
contradiction to strict positivity of the maximum unless
$\sup(U-V^\delta)\le0$.

Therefore $U\le V^\delta$ for every $\delta>0$. Letting
$\delta\downarrow0$ proves comparison. Existence follows by solving the
compactly truncated HJBI on $K_R$, extending with stopped dynamics, and using
the Lyapunov exit estimate from the previous theorem to pass to the limit.
Comparison identifies every subsequential limit satisfying the same weighted
boundary condition at infinity, so the convergence is locally uniform and the
limit is the unique viscosity solution in that fixed-boundary subclass.
\end{proof}

\begin{remark}
This theorem closes the untruncated HJBI comparison bridge for the model class
in this appendix: common Lyapunov subcriticality with $p>1$ tail control,
linear-growth rewards, compact controls/ambiguity, and a fixed weighted
boundary behavior at infinity. It still does not prove differentiability of
boundary or no-quote active sets, nor does it cover state-dependent, power-law,
or empirically learned Hawkes kernels without rechecking the same weighted
comparison assumptions.
\end{remark}

\section{Interior Quote Optimizer Sensitivity}

The HJBI theorems above identify the value for the stated control class. This
section records the additional local condition needed to connect the robust
Hamiltonian to differentiable quote rules as the branching matrix or ambiguity
parameter changes. It covers the smooth interior region of the implemented
quote map; no-quote boundaries and active-set switches are intentionally
handled as nonsmooth cases.

Fix a state/adjoint tuple $z$ collecting $(t,q,\lambda)$ and the value
increments entering the local HJBI Hamiltonian. Let $\theta\in\Theta\subset
\R^m$ be a finite-dimensional parameterization of either $\Gamma$ itself or an
ambiguity set centered at $\Gamma$. For a quote vector $\delta\in D\subset
\R^k$, define the reduced robust Hamiltonian
\[
\mathfrak H(\delta,\theta;z)
=\inf_{\Gamma'\in\mathcal A(\theta)}
\mathcal H(\delta,\Gamma';z),
\]
where $\mathcal H$ is the local reward-plus-generator expression in the HJBI.

\begin{assumption}[Smooth interior robust optimizer]
For the fixed $z$, assume there are neighborhoods of
$(\delta_0,\theta_0)$ such that:
\begin{enumerate}
\item $D$ is open around $\delta_0$, and the admissible quote constraints are
inactive in that neighborhood.
\item Nature has a unique active minimizer
$\Gamma^\dagger(\delta,\theta;z)$, and this selector is $C^1$ in
$(\delta,\theta)$.
\item The reduced Hamiltonian $\mathfrak H$ is $C^2$ in $\delta$ and has
continuous mixed derivative $D^2_{\delta\theta}\mathfrak H$.
\item The quote objective is uniformly strictly concave:
\[
D^2_{\delta\delta}\mathfrak H(\delta,\theta;z)\preceq -m I_k
\]
for some $m>0$ throughout the neighborhood.
\item $\delta_0$ is the unique local maximizer at $\theta_0$, so
$D_\delta\mathfrak H(\delta_0,\theta_0;z)=0$.
\end{enumerate}
\end{assumption}

\begin{theorem}[Interior optimizer differentiability]
Under the smooth interior robust optimizer assumption, there exists a
neighborhood $N$ of $\theta_0$ and a unique interior maximizer
$\delta^*(\theta;z)$ in the corresponding quote neighborhood. The map
$\theta\mapsto\delta^*(\theta;z)$ is $C^1$ on $N$, and
\[
D_\theta\delta^*(\theta_0;z)
=
-\left[
D^2_{\delta\delta}\mathfrak H(\delta_0,\theta_0;z)
\right]^{-1}
D^2_{\delta\theta}\mathfrak H(\delta_0,\theta_0;z).
\]
In the one-dimensional half-spread case, if $e$ is a Perron-adverse structural
direction and
\[
D^2_{\delta\theta}\mathfrak H(\delta_0,\theta_0;z)e\ge0,
\]
then $D_\theta\delta^*(\theta_0;z)e\ge0$; the robust interior half-spread
widens to first order in that adverse direction.
\end{theorem}

\begin{proof}
By the smooth active-minimizer assumption and the envelope/Danskin theorem,
the reduced Hamiltonian has the stated derivatives in a neighborhood of
$(\delta_0,\theta_0)$. Define
\[
G(\delta,\theta)=D_\delta\mathfrak H(\delta,\theta;z).
\]
The first-order condition at the reference point is $G(\delta_0,\theta_0)=0$.
The derivative of $G$ with respect to $\delta$ is
$D^2_{\delta\delta}\mathfrak H(\delta_0,\theta_0;z)$, which is nonsingular
because it is negative definite. The implicit function theorem therefore gives
a unique $C^1$ function $\delta^*(\theta;z)$ solving
$G(\delta^*(\theta;z),\theta)=0$ locally. Strict concavity makes that critical
point the unique local maximizer. Differentiating the identity
$G(\delta^*(\theta;z),\theta)=0$ gives the displayed formula. In dimension
one, the denominator $D^2_{\delta\delta}\mathfrak H$ is negative, so the sign
of the directional derivative agrees with the sign of
$D^2_{\delta\theta}\mathfrak H e$. This proves the widening statement.
\end{proof}

\begin{remark}
The theorem is local by design. If the optimal quote hits a no-quote boundary,
if a cap such as $\delta\le\delta_{\max}$ binds, or if two worst-case Hawkes
matrices exchange activity, the optimizer may have only one-sided directional
derivatives or a kink. Those cases are consistent with the implemented
liquidity-guard policy, but they should be reported as nonsmooth active-set
behavior rather than as classical differentiability.
\end{remark}

\section{Consolidated Multitype Control Theorem}

The preceding sections are intentionally modular because compact verification,
untruncated stability, comparison, and quote sensitivity require different
assumptions. The theorem below states the end-to-end control result that the
manuscript may cite.

\begin{theorem}[End-to-end multitype Hawkes HJBI/control result]
Consider the finite-horizon robust market-making game with finite inventory
grid, exponential $d$-type Hawkes intensity state $\lambda\in\R_+^d$, compact
quote/no-quote control graph, compact rectangular ambiguity set $\mathcal A$,
controlled mark rates bounded by the Hawkes intensity state, and continuous
rewards and terminal payoff with at most linear $W_v$-growth. Suppose common
Lyapunov subcriticality with the $p>1$ tail-integrability condition and
weighted comparison regularity hold, and compare candidate viscosity solutions
only in a $W_v$-weighted subclass with the same stated boundary condition at
infinity. Then:
\begin{enumerate}
\item the controlled Hawkes PDMP is nonexplosive for every admissible
controller and Nature selector;
\item the robust finite-horizon value is finite in the $W_v$-weighted class;
\item the value is the unique viscosity solution, in that weighted class and
with the fixed stated boundary condition at infinity, of the untruncated
multitype Hawkes HJBI equation;
\item compact Lyapunov-sublevel truncations converge locally uniformly to the
untruncated value as the truncation radius tends to infinity;
\item if the untruncated value is classically differentiable, has the same
$W_v$-growth, its generator is integrable under the selected localized
dynamics, and the Isaacs saddle selectors exist with admissible localization
limits, then those selectors verify the untruncated game value; and
\item on smooth interior quote regions satisfying the smooth interior robust
optimizer assumption--inactive quote constraints, a unique $C^1$ active
worst-case structural selector, a $C^2$ reduced Hamiltonian with continuous
mixed derivative, uniformly negative quote Hessian, and a unique local
maximizer--the robust quote map is $C^1$ in the branching/ambiguity parameter
and its derivative is given by the implicit-function formula above.
\end{enumerate}
\end{theorem}

\begin{proof}
Items 1 and 2 are exactly the Lyapunov stability theorem: the common
subcritical vector gives the Foster--Lyapunov drift bound, nonexplosion,
$p>1$ finite weighted maximal moments, and finite linearly growing payoffs.
Item 3 follows from the weighted comparison theorem, which proves comparison for
upper-semicontinuous subsolutions and lower-semicontinuous supersolutions in
the $W_v$-weighted class with the boundary condition at infinity; existence is
obtained by the compact truncation construction. The boundary condition at
infinity is a condition on the compared candidate solutions, not a consequence
of pointwise subcriticality alone. Item 4 is the local-uniform convergence
conclusion of the same comparison-plus-truncation argument, using stopped
Lyapunov-sublevel games and the Lyapunov exit estimate to control paths that
leave compact sublevel sets; it does not identify projected-boundary controls
with stopped-boundary controls without this localization error bound. Item 5
is the classical verification argument localized at $\tau_R$, followed by
Dynkin's formula, the max--min inequalities, the weighted moment bounds, and
then $R\uparrow\infty$ under the stated generator integrability and selector
localization assumptions. Item 6 is exactly the interior optimizer
differentiability theorem, with its smooth active-minimizer and strict
concavity hypotheses included explicitly in the theorem statement.
\end{proof}

\section{Bivariate Perron Bridge}

This section records the bivariate bridge theorem used by the manuscript. It is
still a reduced-form risk-premium theorem, not a complete HJBI proof, but it
upgrades the scalar statement to the smallest genuinely multitype Hawkes
setting.

Let
\[
\Gamma =
\begin{pmatrix}
\gamma_{11} & \gamma_{12}\\
\gamma_{21} & \gamma_{22}
\end{pmatrix}
\]
be nonnegative, irreducible, and subcritical with Perron root
$\rho=\rho(\Gamma)<1$. Let $r,\ell>0$ be right/left Perron vectors normalized
by $\ell^\top r=1$. Let the second eigenvalue be $\eta$, and assume a fixed
gap condition along a near-critical sequence,
\[
\rho\uparrow1,\qquad |\rho-\eta|\ge g_0>0.
\]
For two-type order flow with baseline $\mu\in\R^2_+$, the stationary intensity
is
\[
\bar\lambda=(I-\Gamma)^{-1}\mu.
\]
For a reduced market-making risk direction $w\in\R^2_+$, define
\[
V_{\mathrm{bi}}(\Gamma)
 =
C\left(w^\top (I-\Gamma)^{-1}\mu\right)^2,
\qquad C>0.
\]
We work on a compact Perron-regular class with uniformly bounded eigenprojection
conditioning and Perron visibility bounded below:
\[
A:=(w^\top r)(\ell^\top\mu)\ge a_0>0.
\]

\begin{lemma}[Bivariate resolvent decomposition]
Under the bivariate assumptions above,
\[
(I-\Gamma)^{-1}
 =
\frac{r\ell^\top}{1-\rho}
+R_\eta,
\qquad
\|R_\eta\|\le C_\eta,
\]
where $C_\eta$ depends only on the spectral gap, vector normalization,
eigenprojection conditioning, and a compact upper bound on $\|\Gamma\|$.
\end{lemma}

\begin{proof}
Because $\Gamma$ is a real $2\times2$ matrix with simple Perron root and
separated second eigenvalue, the spectral projection onto the Perron eigenspace
is $P=r\ell^\top$. The complementary projection is $I-P$. On the Perron
subspace, $(I-\Gamma)^{-1}$ has eigenvalue $(1-\rho)^{-1}$. On the complementary
one-dimensional subspace, the resolvent has eigenvalue $(1-\eta)^{-1}$ times a
bounded projection. Since $|\rho-\eta|\ge g_0$ and $\rho<1$, the complementary
term is uniformly bounded on the stated compact class. This gives the
decomposition.
\end{proof}

\begin{lemma}[Bivariate risk singularity]
Under the same assumptions,
\[
V_{\mathrm{bi}}(\Gamma)
 =
C A^2(1-\rho)^{-2}
 +O((1-\rho)^{-1}).
\]
Consequently $V_{\mathrm{bi}}(\Gamma)=\Theta((1-\rho)^{-2})$ as
$\rho\uparrow1$.
\end{lemma}

\begin{proof}
By the resolvent decomposition,
\[
w^\top(I-\Gamma)^{-1}\mu
=
\frac{(w^\top r)(\ell^\top\mu)}{1-\rho}
+ w^\top R_\eta \mu
=
\frac{A}{1-\rho}+O(1).
\]
Squaring yields
\[
\left(\frac{A}{1-\rho}+O(1)\right)^2
=
A^2(1-\rho)^{-2}+O((1-\rho)^{-1})+O(1).
\]
Since $A>0$, the leading term is nonzero and proves the stated singularity.
\end{proof}

\begin{assumption}[Perron-coefficient regularity]
For the local perturbation paths below, the Perron visibility coefficient
\[
A(\Gamma):=(w^\top r(\Gamma))(\ell(\Gamma)^\top\mu)
\]
does not cancel the gap singularity. Concretely, for relative-slack paths we
assume $A(\Gamma+E_{\mathrm{rel}})=A(\Gamma)+o(\eps_{\mathrm{rel}})$, and for
absolute paths we assume $A(\Gamma+E_{\mathrm{abs}})=A(\Gamma)+o(\eps)$. A
sufficient special case is a spectral path that changes the Perron root while
holding the Perron eigendirections fixed to first order.
\end{assumption}

\begin{theorem}[Perron-aligned bivariate relative-slack premium]
Let $E_{\mathrm{rel}}$ be a perturbation for which
\[
\rho(\Gamma+E_{\mathrm{rel}})
=
\rho+\eps_{\mathrm{rel}}(1-\rho)+O(\eps_{\mathrm{rel}}^2(1-\rho)),
\]
and suppose $\Gamma+E_{\mathrm{rel}}$ remains in the same compact bivariate
Perron-regular class with Perron visibility bounded below and satisfying
Perron-coefficient regularity. Then
\[
V_{\mathrm{bi}}(\Gamma+E_{\mathrm{rel}})
-V_{\mathrm{bi}}(\Gamma)
=
2CA^2\eps_{\mathrm{rel}}(1-\rho)^{-2}
+O(\eps_{\mathrm{rel}}^2(1-\rho)^{-2})
+o(\eps_{\mathrm{rel}}(1-\rho)^{-2})
+O(\eps_{\mathrm{rel}}(1-\rho)^{-1}).
\]
In particular, the leading singular order is
$\Theta(\eps_{\mathrm{rel}}(1-\rho)^{-2})$. Under local quote sensitivity, a
finite ambiguity set that contains and activates this Perron-aligned adversary
has the same square-law half-spread order.
\end{theorem}

\begin{proof}
The bivariate risk singularity gives
\[
V_{\mathrm{bi}}(\Gamma)=C A^2(1-\rho)^{-2}+O((1-\rho)^{-1}).
\]
The perturbed gap satisfies
\[
1-\rho(\Gamma+E_{\mathrm{rel}})
=(1-\eps_{\mathrm{rel}})(1-\rho)+O(\eps_{\mathrm{rel}}^2(1-\rho)).
\]
Uniform Perron visibility and compactness keep the leading coefficient bounded
above and below. Therefore the leading change is the same calculation as in the
scalar relative-slack theorem with $a=2$:
\[
C A^2(1-\rho)^{-2}\{(1-\eps_{\mathrm{rel}})^{-2}-1\}
=
2 C A^2\eps_{\mathrm{rel}}(1-\rho)^{-2}
+O(\eps_{\mathrm{rel}}^2(1-\rho)^{-2}).
\]
The lower bound follows because $A$ is bounded away from zero and the
Perron-coefficient regularity prevents first-order cancellation; the
complementary resolvent term is lower order. Local quote sensitivity transfers
the order to the half-spread correction.
\end{proof}

\begin{theorem}[Perron-aligned bivariate absolute-ambiguity derivative law]
Let $E_{\mathrm{abs}}$ satisfy
\[
\rho(\Gamma+E_{\mathrm{abs}})
=\rho+\eps+O(\eps^2),
\qquad \eps=o(1-\rho),
\]
and suppose the perturbed matrices remain in the same bivariate Perron-regular
class with Perron visibility bounded below and satisfying Perron-coefficient
regularity. Then
\[
V_{\mathrm{bi}}(\Gamma+E_{\mathrm{abs}})
-V_{\mathrm{bi}}(\Gamma)
=
2CA^2\eps(1-\rho)^{-3}
+O(\eps^2(1-\rho)^{-4})
+o(\eps(1-\rho)^{-3})
+O(\eps(1-\rho)^{-2}).
\]
In particular, the leading singular order is
$\Theta(\eps(1-\rho)^{-3})$. Under local quote sensitivity, a finite ambiguity
set that contains and activates this Perron-aligned adversary has the same
derivative-law half-spread order.
\end{theorem}

\begin{proof}
Using the same leading expansion,
\[
V_{\mathrm{bi}}(\Gamma+\!E_{\mathrm{abs}})-V_{\mathrm{bi}}(\Gamma)
=
C A^2\left((1-\rho-\eps+O(\eps^2))^{-2}
-(1-\rho)^{-2}\right)
+O(\eps(1-\rho)^{-2}).
\]
Since $\eps=o(1-\rho)$, a Taylor expansion of $x^{-2}$ at $x=1-\rho$ gives
\[
(1-\rho-\eps)^{-2}-(1-\rho)^{-2}
=
2\eps(1-\rho)^{-3}+O(\eps^2(1-\rho)^{-4}).
\]
The $O(\eps(1-\rho)^{-2})$ contribution from the bounded complementary resolvent
and the $o(\eps(1-\rho)^{-3})$ contribution from the Perron coefficient are
lower order. Positivity of the aligned Perron coefficient gives the matching
lower bound. Local quote sensitivity again transfers the order.
\end{proof}

\begin{remark}[How to realize the perturbations]
The first-order Perron perturbation proposition gives a concrete construction:
take $E$ proportional to the Frobenius-normalized Perron direction
$\ell r^\top$. Then
\[
\rho(\Gamma+tE)-\rho(\Gamma)=t\langle E,\ell r^\top\rangle_F+O(t^2),
\]
so choosing $t=O(\eps_{\mathrm{rel}}(1-\rho))$ gives the relative theorem and
choosing $t=O(\eps)$ gives the absolute theorem. This is exactly the adversarial
direction used in the matrix-resolvent experiments.
\end{remark}

\section{Conclusion}

The square-law statement is correct only after defining the ambiguity radius
relative to the remaining critical slack or to an already-normalized risk
quantity. If $\eps$ is an absolute branching-matrix radius, the robust premium
inherits one additional derivative of the critical singularity. The same
distinction holds in the bivariate Perron-visible setting. This appendix
therefore converts the original MESA claim into precise, testable scalar and
bivariate theorem statements.

\end{document}
